\chapter{Widget catalogue: layout primitives}
\label{ch:catalogue_layout}

\begin{aside}
A reference of every layout primitive in \maya{}'s built-in
catalogue. Quick recipes, not deep dives.
\end{aside}

\section{column / row}

Vertical / horizontal flex container. Default growth: 1
for all children.

\section{stack}

Z-ordered stack. Children paint in order; last wins.

\section{center}

Centres its child both axes. Equivalent to row(spacer,
child, spacer) | column(spacer, ., spacer) but cheaper.

\section{align}

Specify alignment on each axis:
\code{align(left, top)}.

\section{padding}

Reserve cells around the child. \code{padding(1, 2)} =
1 row, 2 columns.

\section{margin}

Like padding but outside the child's box.

\section{border}

Wrap with a border preset. See Chapter~\ref{ch:borders}.

\section{spacer}

A flexible empty element that pushes siblings.

\section{box}

A fixed-size empty element. Useful for padding when
flex won't do.

\section{flex}

A modifier that sets the flex factor.

\section{width / height}

Constraints. \code{width(20)} = exactly 20 cells.
\code{width(20, 80)} = between 20 and 80.

\section{overflow}

Wrap, ellipsis, scroll, truncate, hidden.

\section{clip}

Force a clip rectangle. Children paint within.

\section{Recap}

\begin{recap}
\begin{itemize}[leftmargin=*]
  \item Containers: column, row, stack.
  \item Positioning: center, align, padding, margin.
  \item Decoration: border, spacer, box.
  \item Sizing: flex, width, height.
  \item Overflow and clip controls.
\end{itemize}
\end{recap}

\section*{Workshop}

\begin{enumerate}[leftmargin=*]
  \item Build a layout using only catalogue primitives.
  \item Replace a custom widget with composed primitives.
  \item Time the difference between flex(1) and explicit
    sizes.
\end{enumerate}
